% Source: https://artofproblemsolving.com/community/c3h3072883_some_difficult_counting_problems
% Stable identity: aops-topic-3072883-q06
The graph is connected: a component with $k$ vertices would require $k-1\ge\dfrac{n}{2}$, so two components would together contain more than $n$ vertices. Choose a longest simple path $v_1,\ldots,v_k$. Every neighbor of either endpoint lies on the path, or the path could be extended. Among indices $i=1,\ldots,k-1$, at least one has both edges $v_1v_{i+1}$ and $v_iv_k$: otherwise the two sets of such indices would be disjoint, implying $\deg(v_1)+\deg(v_k)\le k-1$, contradicting their sum at least $n\ge k$. These edges close a cycle through all path vertices, traversed as $v_1,v_2,\ldots,v_i,v_k,v_{k-1},\ldots,v_{i+1},v_1$. If $k<n$, connectedness supplies an edge from a vertex outside this cycle to a vertex on it. Starting at that outside vertex and going around the cycle produces a path with $k+1$ vertices, contradicting maximality. Thus $k=n$ and the cycle is Hamiltonian.
\begin{asy}
usepackage("amsmath");
size(7cm);
pair[] V;for(int i=0;i<6;++i)V.push(dir(90-60i));for(int i=0;i<6;++i){draw(V[i]--V[(i+1)%6],blue);dot(V[i]);label("$v_{"+string(i+1)+"}$",V[i],unit(V[i]));}draw(V[0]--V[3],dashed);draw(V[1]--V[4],dashed);draw(V[2]--V[5],dashed);label("example with degree 3",(0,-1.5));
\end{asy}
